%=====================================================================
\chapter{Concept Learning and the Version Space}\label{ch:concept}
%=====================================================================
\locator{1}{Part I --- Foundations}

\begin{outcomes}
By the end of this chapter you will be able to:
\begin{tight}
  \item \textbf{Represent} a concept as a conjunction of attribute constraints,
        and \textbf{order} hypotheses from specific to general.
  \item \textbf{Trace} the \textsc{Find-S} algorithm on a set of training
        examples, and \textbf{state} its limitations.
  \item \textbf{Define} the version space, and \textbf{trace} the
        \textsc{Candidate-Elimination} algorithm through the $S$ and $G$
        boundaries by hand.
  \item \textbf{Classify} a new instance with a partially learned concept, and
        \textbf{choose} the most informative query.
  \item \textbf{Explain} why an unbiased learner cannot generalise, and
        \textbf{state} the inductive bias of \textsc{Candidate-Elimination}.
\end{tight}
\end{outcomes}

\begin{prereq}
\chref{introduction} (hypotheses, hypothesis spaces and inductive bias,
introduced informally in \dsref{hypothesis-space}). This chapter makes those
ideas exact.
\end{prereq}

\begin{keyterms}
concept learning $\bullet$ instance space $\bullet$ target concept $\bullet$
attribute constraint $\bullet$ consistent hypothesis $\bullet$
more-general-than-or-equal-to $\bullet$ \textsc{Find-S} $\bullet$ version space
$\bullet$ \textsc{List-Then-Eliminate} $\bullet$ general boundary $G$ $\bullet$
specific boundary $S$ $\bullet$ \textsc{Candidate-Elimination} $\bullet$ query
$\bullet$ unbiased learner $\bullet$ inductive bias
\end{keyterms}

\section{Concept Learning as a Task}

Much of learning is learning a \emph{category}: which releases are ready to ship,
which tasks will slip, which tickets are urgent. Each such category is a
\term{concept}, and it can be described as a yes/no function over the things
being classified.

\begin{definitionbox}[Concept Learning]
\term{Concept learning} is inferring a boolean-valued function from training
examples of its input and output. Formally:
\begin{tight}
  \item the \term{instance space} $X$ is the set of all things that could be
        classified;
  \item the \term{target concept} is an unknown function $c: X \to \{0, 1\}$;
  \item the training examples $D$ are pairs $\langle x, c(x)\rangle$;
  \item the learner searches a hypothesis space $H$ for an $h$ with $h(x) = c(x)$
        for every $x \in X$.
\end{tight}
\end{definitionbox}

The running example of this chapter is a \emph{change advisory board}, which
decides whether each software release may go to production. Each release is
described by six attributes, and four of the board's decisions have been
observed. (The table is Mitchell's \emph{EnjoySport} example with its attributes
renamed, so every number in this chapter can be checked against the set text.)

\rowcolors{2}{white}{lightbg}
\begin{longtable}{@{}L{1.4cm}L{1.3cm}L{1.6cm}L{1.9cm}L{1.8cm}L{1.6cm}L{1.6cm}L{1.9cm}@{}}
\toprule
\rowcolor{primary!15}
\textbf{Release} & \textbf{Tests} & \textbf{Review} & \textbf{Docs} & \textbf{Rollback} &
\textbf{Traffic} & \textbf{Scope} & \textbf{Approve} \\
\midrule
\endhead
1 & Pass & Done & Complete & Ready & Quiet & Same & Yes \\
2 & Pass & Done & Partial & Ready & Quiet & Same & Yes \\
3 & Fail & Pending & Partial & Ready & Quiet & Changed & No \\
4 & Pass & Done & Partial & Ready & Peak & Changed & Yes \\
\bottomrule
\end{longtable}

\noindent Tests is Pass, Flaky or Fail; Review is Done or Pending; Docs is
Complete or Partial; Rollback (a tested plan to undo the release) is Ready or
None; Traffic, the load expected at release time, is Quiet or Peak; and Scope
records whether the scope has Changed since the last review.

\begin{definitionbox}[Hypotheses as Conjunctions of Constraints]
Each hypothesis is a vector of six \term{constraints}, one per attribute. Each
constraint is one of:
\begin{tight}
  \item a specific value, such as \emph{Done} --- only that value is acceptable;
  \item ``?'' --- any value is acceptable;
  \item ``$\varnothing$'' --- no value is acceptable.
\end{tight}
A hypothesis classifies an instance as positive if the instance satisfies every
constraint. So $\langle$Pass, ?, ?, Ready, ?, ?$\rangle$ means ``tests pass and a
rollback is ready, whatever else''; $\langle$?, ?, ?, ?, ?, ?$\rangle$ approves
every release; and a hypothesis containing $\varnothing$ approves none.
\end{definitionbox}

\begin{worked}[How Big Are These Spaces?]
Tests has 3 values and the other five attributes have 2 each.
\begin{tight}
  \item \textbf{Instances:} $3 \times 2^5 = \mathbf{96}$ possible releases.
  \item \textbf{Hypotheses, as written:} each constraint can also be ``?'' or
        ``$\varnothing$'', so $5 \times 4^5 = 5120$.
  \item \textbf{Hypotheses, by meaning:} every hypothesis containing a
        $\varnothing$ classifies every release as negative, so they all mean the
        same thing. Counting that as one: $1 + 4 \times 3^5 = \mathbf{973}$.
  \item \textbf{All possible concepts} over 96 releases: every subset of the
        releases is a concept, so $2^{96} \approx 7.9 \times 10^{28}$.
\end{tight}
\smallskip
\textbf{The conjunctive language can express 973 of the $7.9 \times 10^{28}$
possible concepts.} That restriction looks like a weakness. By the end of the
chapter it will turn out to be the only reason learning is possible at all.
\end{worked}

\section{The Inductive Learning Hypothesis}

A learner sees only the training examples. The best it can guarantee is a
hypothesis that fits them. Everything rests on an assumption, which Mitchell
names explicitly:

\begin{conceptbox}[The Inductive Learning Hypothesis]
Any hypothesis found to approximate the target function well over a sufficiently
large set of training examples will also approximate the target function well
over other unobserved examples.\\[3pt]
This is an assumption, not a theorem. It fails whenever the future does not
resemble the training data --- which is why \chref{generalization} and
\chref{evaluation} measure performance on held-out data rather than trusting it.
\end{conceptbox}

\section{The General-to-Specific Ordering}

Hypotheses are not an unordered heap. Some are strictly more permissive than
others, and that structure is what makes efficient search possible.

\begin{definitionbox}[More-General-Than-or-Equal-To]
Hypothesis $h_j$ is \term{more general than or equal to} $h_k$, written
$h_j \ge_g h_k$, if every instance that $h_k$ classifies as positive, $h_j$ also
classifies as positive. $h_j$ is strictly more general ($h_j >_g h_k$) if, in
addition, $h_k \not\ge_g h_j$. The relation is a \term{partial order}: some
pairs of hypotheses are not comparable at all.
\end{definitionbox}

\dsfig{%
\begin{tikzpicture}[font=\scriptsize,
  hn/.style={rounded corners=2pt, draw=primary, fill=primary!8, text=primary!70!black,
             font=\scriptsize, minimum height=0.55cm, inner xsep=4pt}]
% instance space
\draw[rounded corners=6pt, draw=gray!50, fill=gray!4] (0,0) rectangle (6.2,4.2);
\node[font=\small\bfseries, text=gray!55!black, anchor=north west] at (0.1,4.15)
     {Instances $X$};
\draw[draw=secondary, line width=0.9pt, fill=secondary!8] (3.1,1.95) ellipse (2.65cm and 1.6cm);
\draw[draw=greenhl!70!black, line width=0.9pt, fill=greenhl!12] (2.35,1.75) ellipse (1.3cm and 0.85cm);
\draw[draw=goldhl, line width=0.9pt, fill=goldhl!14, fill opacity=0.7] (3.95,1.75) ellipse (1.3cm and 0.85cm);
\node[font=\tiny\bfseries, text=secondary!70!black] at (3.1,3.2) {covered by $h_2$};
\node[font=\tiny\bfseries, text=greenhl!45!black] at (1.75,1.75) {$h_1$};
\node[font=\tiny\bfseries, text=goldhl!45!black] at (4.55,1.75) {$h_3$};
\fill[accent] (3.15,1.85) circle (2pt) node[below, font=\tiny, text=accent] {$x_1$};
\fill[accent] (4.9,2.75) circle (2pt) node[right, font=\tiny, text=accent] {$x_2$};
% hypothesis space
\draw[rounded corners=6pt, draw=gray!50, fill=gray!4] (7.0,0) rectangle (16.6,4.2);
\node[font=\small\bfseries, text=gray!55!black, anchor=north west] at (7.1,4.15)
     {Hypotheses $H$};
\node[hn, draw=secondary, fill=secondary!10, text=secondary!70!black] (h2) at (11.4,3.1)
     {$h_2 = \langle$Pass, ?, ?, ?, ?, ?$\rangle$};
\node[hn, draw=greenhl!70!black, fill=greenhl!12, text=greenhl!40!black, align=center]
     (h1) at (8.85,1.2) {$h_1 =$\\$\langle$Pass, ?, ?, Ready, ?, ?$\rangle$};
\node[hn, draw=goldhl, fill=goldhl!14, text=goldhl!40!black, align=center]
     (h3) at (13.95,1.2) {$h_3 =$\\$\langle$Pass, ?, ?, ?, Peak, ?$\rangle$};
\draw[-{Stealth[length=2mm]}, gray!60, line width=0.9pt] (h1) -- (h2);
\draw[-{Stealth[length=2mm]}, gray!60, line width=0.9pt] (h3) -- (h2);
\node[font=\tiny, text=gray!55!black, align=center] at (11.4,1.2)
     {$h_1$, $h_3$:\\not\\comparable};
\draw[-{Stealth[length=2.5mm]}, primary!50, line width=1.2pt] (16.25,0.3) -- (16.25,3.9);
\node[font=\tiny, text=primary, rotate=90, anchor=south] at (16.2,2.1) {more general};
\node[anchor=north west, align=left, text width=16.4cm, font=\scriptsize, text=gray!55!black]
     at (0,-0.15)
     {$x_1 = \langle$Pass, Done, Partial, Ready, Peak, Same$\rangle$ satisfies all three;
      $x_2 = \langle$Pass, Done, Partial, None, Quiet, Same$\rangle$ satisfies only $h_2$.
      $h_2 \ge_g h_1$ and $h_2 \ge_g h_3$, because the set each covers is inside the set
      $h_2$ covers.};
\end{tikzpicture}}%
{The more-general-than ordering. A more general hypothesis covers a superset of the
instances; hypotheses whose covered sets merely overlap, like $h_1$ and $h_3$, are not
comparable.}{general-specific}

\section{Find-S: The Most Specific Consistent Hypothesis}

The simplest algorithm that uses the ordering starts at the most specific
hypothesis and generalises only as far as each positive example forces it to.

\begin{definitionbox}[The Find-S Algorithm]
\begin{tightnum}
  \item Initialise $h$ to the most specific hypothesis
        $\langle\varnothing, \varnothing, \varnothing, \varnothing, \varnothing,
        \varnothing\rangle$.
  \item For each \emph{positive} training example $x$: for each constraint in
        $h$ that $x$ does not satisfy, replace it by the next more general
        constraint that $x$ does satisfy ($\varnothing \to$ the value of $x$;
        a different value $\to$ ``?'').
  \item Ignore negative examples. Output $h$.
\end{tightnum}
\end{definitionbox}

\begin{worked}[Find-S on the Release Decisions]
\textbf{Start:} $h_0 = \langle\varnothing, \varnothing, \varnothing, \varnothing,
\varnothing, \varnothing\rangle$.

\textbf{Release 1 (+):} nothing is satisfied, so copy the example:
$h_1 = \langle$Pass, Done, Complete, Ready, Quiet, Same$\rangle$.

\textbf{Release 2 (+):} Docs is Partial, not Complete, so generalise that
constraint: $h_2 = \langle$Pass, Done, ?, Ready, Quiet, Same$\rangle$.

\textbf{Release 3 ($-$):} ignored. (It is correctly classified as negative anyway
--- its tests failed.)

\textbf{Release 4 (+):} Traffic and Scope differ, so generalise both:
$h_4 = \langle$Pass, Done, ?, Ready, ?, ?$\rangle$.

\smallskip
\textbf{Output:} ``tests pass, review done and a rollback ready''. Why is it safe
to ignore negatives? If the target concept is in $H$, it is at least as general as
$h$ (which only ever generalised as far as positives forced). A negative example
is excluded by the target, so it is excluded by the more specific $h$ as well.
\end{worked}

\begin{alertbox}[What Find-S Cannot Tell You]
\textsc{Find-S} returns \emph{a} consistent hypothesis, but not whether it is the
only one; it cannot say whether it has converged; it cannot detect noisy or
inconsistent data, because it never looks at negatives; and it always prefers the
most specific answer, for no reason other than that it started there. The version
space fixes the first three.
\end{alertbox}

\section{Version Spaces}

\begin{definitionbox}[Consistent Hypothesis and Version Space]
A hypothesis $h$ is \term{consistent} with training examples $D$ if $h(x) = c(x)$
for every example $\langle x, c(x)\rangle$ in $D$. The \term{version space}
$VS_{H,D}$ is the set of all hypotheses in $H$ consistent with $D$ --- every
hypothesis the data has not yet ruled out.
\end{definitionbox}

The obvious way to find it is \textsc{List-Then-Eliminate}: write down every
hypothesis in $H$, and for each training example cross out the ones it
contradicts. What remains is the version space. For this example's 973 hypotheses
this is feasible; for any realistic hypothesis space it is not. The trick of the
next section is that the version space never needs to be listed --- two
boundaries are enough.

\section{The Candidate-Elimination Algorithm}

\begin{definitionbox}[General and Specific Boundaries]
The \term{general boundary} $G$ is the set of maximally general hypotheses in the
version space; the \term{specific boundary} $S$ is the set of maximally specific
ones. Every hypothesis in the version space lies between them: it is more general
than or equal to some member of $S$ and more specific than or equal to some member
of $G$.
\end{definitionbox}

\begin{definitionbox}[The Candidate-Elimination Algorithm]
Initialise $G$ to $\{\langle ?, ?, ?, ?, ?, ?\rangle\}$ and $S$ to
$\{\langle\varnothing, \varnothing, \varnothing, \varnothing, \varnothing,
\varnothing\rangle\}$. For each training example $d$:
\begin{tight}
  \item \textbf{If $d$ is positive:} remove from $G$ every hypothesis that does not
        cover $d$. Replace each $s \in S$ that does not cover $d$ by its minimal
        generalisations that do, keeping only those that are more specific than
        some member of $G$.
  \item \textbf{If $d$ is negative:} remove from $S$ every hypothesis that covers
        $d$. Replace each $g \in G$ that covers $d$ by its minimal
        specialisations that do not, keeping only those that are more general
        than some member of $S$, and then drop any that is less general than
        another member of $G$.
\end{tight}
\end{definitionbox}

\dsfig{%
\begin{tikzpicture}[font=\scriptsize]
\fill[primary!4] (0,0) rectangle (13.6,3.6);
\draw[secondary, line width=1.6pt] (0.6,2.9) .. controls (4,2.6) and (9,2.6) .. (13.0,2.9);
\draw[greenhl!70!black, line width=1.6pt] (0.6,0.7) .. controls (4,1.0) and (9,1.0) .. (13.0,0.7);
\fill[purplehl!10] (0.6,2.9) .. controls (4,2.6) and (9,2.6) .. (13.0,2.9)
     -- (13.0,0.7) .. controls (9,1.0) and (4,1.0) .. (0.6,0.7) -- cycle;
\draw[secondary, line width=1.6pt] (0.6,2.9) .. controls (4,2.6) and (9,2.6) .. (13.0,2.9);
\draw[greenhl!70!black, line width=1.6pt] (0.6,0.7) .. controls (4,1.0) and (9,1.0) .. (13.0,0.7);
\node[font=\scriptsize\bfseries, text=secondary!70!black, anchor=south east] at (13.0,2.95)
     {$G$ --- the most general hypotheses still consistent};
\node[font=\scriptsize\bfseries, text=greenhl!45!black, anchor=north west] at (0.6,0.62)
     {$S$ --- the most specific hypotheses still consistent};
\node[font=\small\bfseries, text=purplehl!70!black] at (6.8,1.8) {version space};
\draw[-{Stealth[length=2.5mm]}, accent, line width=1.2pt] (2.2,3.35) -- (2.2,2.55);
\node[font=\tiny, text=accent, anchor=west, align=left] at (2.35,3.3)
     {negative examples\\push $G$ down (specialise)};
\draw[-{Stealth[length=2.5mm]}, greenhl!60!black, line width=1.2pt] (11.4,0.2) -- (11.4,1.05);
\node[font=\tiny, text=greenhl!45!black, anchor=east, align=right] at (11.25,0.25)
     {positive examples\\push $S$ up (generalise)};
\node[font=\tiny, text=gray!55!black, rotate=90] at (-0.3,1.8) {specific \quad$\longrightarrow$\quad general};
\end{tikzpicture}}%
{The version space is everything between the two boundaries. Each example moves one
boundary towards the other; if they meet, the concept has been learned exactly.}{boundaries}

\begin{worked}[Candidate-Elimination on the Release Decisions]
\textbf{Start.} $S_0 = \{\langle\varnothing, \varnothing, \varnothing, \varnothing,
\varnothing, \varnothing\rangle\}$, $\;G_0 = \{\langle ?, ?, ?, ?, ?, ?\rangle\}$.

\smallskip
\textbf{Release 1 (+)} $\langle$Pass, Done, Complete, Ready, Quiet, Same$\rangle$.
$G_0$ covers it: keep. $S_0$ does not: generalise minimally.\\
$S_1 = \{\langle$Pass, Done, Complete, Ready, Quiet, Same$\rangle\}$, $\;G_1 = G_0$.

\smallskip
\textbf{Release 2 (+)} $\langle$Pass, Done, Partial, Ready, Quiet, Same$\rangle$.
Docs differs, so it becomes ``?''.\\
$S_2 = \{\langle$Pass, Done, ?, Ready, Quiet, Same$\rangle\}$, $\;G_2 = G_0$.

\smallskip
\textbf{Release 3 ($-$)} $\langle$Fail, Pending, Partial, Ready, Quiet, Changed$\rangle$.
$S_2$ already excludes it. $G_2$ covers it, so specialise $\langle ?, ?, ?, ?, ?,
?\rangle$ one attribute at a time to a value the negative example does not have:
\begin{tight}
  \item Tests: Pass or Flaky. Only \emph{Pass} is still more general than $S_2$.
  \item Review: \emph{Done}. Kept.
  \item Docs: Complete --- but $S_2$ has ``?'' there, so this is not more
        general than $S_2$. Dropped.
  \item Rollback: None --- contradicts $S_2$'s Ready. Dropped.
  \item Traffic: Peak --- contradicts $S_2$'s Quiet. Dropped.
  \item Scope: \emph{Same}. Kept.
\end{tight}
$S_3 = S_2$, $\;G_3 = \{\langle$Pass, ?, ?, ?, ?, ?$\rangle$, $\langle$?, Done, ?,
?, ?, ?$\rangle$, $\langle$?, ?, ?, ?, ?, Same$\rangle\}$.

\smallskip
\textbf{Release 4 (+)} $\langle$Pass, Done, Partial, Ready, Peak, Changed$\rangle$.
$\langle$?, ?, ?, ?, ?, Same$\rangle$ does not cover it: removed from $G$. $S$
generalises Traffic and Scope.\\
$S_4 = \{\langle$Pass, Done, ?, Ready, ?, ?$\rangle\}$, $\;G_4 = \{\langle$Pass, ?,
?, ?, ?, ?$\rangle$, $\langle$?, Done, ?, ?, ?, ?$\rangle\}$.

\smallskip
\textbf{Check.} $S_4$ is exactly what \textsc{Find-S} returned. The difference is
that \textsc{Candidate-Elimination} also knows the other end: the version space
holds six hypotheses, drawn in \dsref{final-vs}.
\end{worked}

\dsfig{%
\begin{tikzpicture}[font=\scriptsize,
  hv/.style={rounded corners=2pt, draw=#1, fill=#1!9, text=#1!55!black,
             font=\scriptsize, minimum height=0.6cm, inner xsep=4pt},
  e/.style={gray!55, line width=0.8pt}]
\node[hv=secondary] (g1) at (2.6,3.2) {$\langle$Pass, ?, ?, ?, ?, ?$\rangle$};
\node[hv=secondary] (g2) at (9.4,3.2) {$\langle$?, Done, ?, ?, ?, ?$\rangle$};
\node[hv=purplehl] (m1) at (0.9,1.6) {$\langle$Pass, ?, ?, Ready, ?, ?$\rangle$};
\node[hv=purplehl] (m2) at (6.0,1.6) {$\langle$Pass, Done, ?, ?, ?, ?$\rangle$};
\node[hv=purplehl] (m3) at (11.1,1.6) {$\langle$?, Done, ?, Ready, ?, ?$\rangle$};
\node[hv=greenhl] (s1) at (6.0,0) {$\langle$Pass, Done, ?, Ready, ?, ?$\rangle$};
\draw[e] (g1) -- (m1); \draw[e] (g1) -- (m2); \draw[e] (g2) -- (m2); \draw[e] (g2) -- (m3);
\draw[e] (m1) -- (s1); \draw[e] (m2) -- (s1); \draw[e] (m3) -- (s1);
\node[font=\scriptsize\bfseries, text=secondary!70!black, anchor=east] at (-1.1,3.2) {$G_4$:};
\node[font=\scriptsize\bfseries, text=greenhl!45!black, anchor=east] at (-1.1,0) {$S_4$:};
\end{tikzpicture}}%
{The final version space for the release decisions: two maximally general hypotheses,
one maximally specific one, and the three between them. A line joins each pair related
by $\ge_g$.}{final-vs}

\section{Using a Partially Learned Concept}

Four examples have not pinned the concept down: six hypotheses survive. They can
still classify new releases, by voting.

\begin{worked}[Classifying by Vote]
\begin{tight}
  \item $\langle$Pass, Done, Complete, Ready, Peak, Changed$\rangle$: satisfies $S_4$,
        and so every hypothesis above it. \textbf{6 yes, 0 no} --- approved with
        the same confidence as if the concept were fully learned.
  \item $\langle$Fail, Pending, Complete, None, Quiet, Same$\rangle$: satisfies
        neither member of $G_4$, so no hypothesis. \textbf{0 yes, 6 no}.
  \item $\langle$Pass, Done, Complete, None, Quiet, Same$\rangle$: \textbf{3 yes,
        3 no}. The data genuinely cannot decide.
  \item $\langle$Pass, Pending, Complete, Ready, Quiet, Same$\rangle$: \textbf{2 yes,
        4 no} --- a majority, but not a certainty.
\end{tight}
\smallskip
Two tests are enough for the unanimous cases: an instance satisfying every
member of $S$ is positive for the whole version space; one satisfying no member
of $G$ is negative for all of it. And if the learner may \emph{ask} the board for a
decision, the best question is the third release: whatever the answer, it
eliminates half the version space. Repeating that halves it each time, so the
concept is found in about $\log_2 |VS|$ queries --- here $\log_2 6 \approx 2.6$,
so three.
\end{worked}

\section{Inductive Bias}

Why stop at conjunctions? Why not use a hypothesis space that can express
\emph{every} one of the $2^{96}$ possible concepts, so that the target is
certainly in it?

\begin{conceptbox}[The Futility of an Unbiased Learner]
Give \textsc{Candidate-Elimination} a hypothesis space containing every subset
of $X$. After any training set, $S$ is the disjunction of the positive examples
seen --- ``release 1 or release 2 or release 4'' --- and $G$ is ``anything except
release 3''. The version space contains every concept between them. For any
release \emph{not} in the training data, exactly half of those concepts call it
positive and half negative: the vote is always a tie. The unbiased learner can
classify the examples it has seen and nothing else. It has memorised; it has not
learned.
\end{conceptbox}

\begin{definitionbox}[Inductive Bias]
The \term{inductive bias} of a learner is the set of assumptions that, added to
the training data, make its classification of new instances follow
\emph{deductively}. The inductive bias of \textsc{Candidate-Elimination} is:
\emph{the target concept is contained in the hypothesis space $H$} --- here, it is
a conjunction of attribute constraints.
\end{definitionbox}

Learners can be ranked by the strength of their bias. A \textbf{rote learner},
which stores examples and classifies only exact matches, has no bias and never
generalises. \textsc{Candidate-Elimination} assumes the target is in $H$, and
classifies a new instance only when the whole version space agrees.
\textsc{Find-S} assumes that as well, and in addition that every instance is
negative unless the evidence forces it to be positive --- a stronger bias, so it
classifies everything. Every algorithm in the rest of this course has a bias,
and its chapter names it.

\begin{alertbox}[The Weakness Every Later Chapter Fixes]
\textsc{Candidate-Elimination} assumes the data is \emph{noise-free}. Record one
approved release as rejected by mistake and the true concept is eliminated; with
enough errors $S$ and $G$ cross, and the version space is empty. Real data always
contains errors. The algorithms of Part III therefore replace ``consistent with
every example'' by ``makes few errors'' and accept a hypothesis that is probably
approximately right. Concept learning is where the vocabulary comes from; it is
not how practical systems are built.
\end{alertbox}

\begin{pitfall}
\begin{tight}
  \item \textbf{Specialising $G$ to values that contradict $S$.} Every new
        member of $G$ must still be more general than a member of $S$; check each
        candidate against $S$ before keeping it.
  \item \textbf{Forgetting to prune $G$ on a positive example.} Members of $G$ that
        do not cover a positive example are removed, as at Release 4.
  \item \textbf{Believing Find-S's answer is unique.} It is one of the
        six hypotheses in the version space --- the most specific.
  \item \textbf{Reading ``unbiased'' as ``better''.} A learner with no bias
        cannot generalise at all.
  \item \textbf{Running Candidate-Elimination on noisy data.} One
        mislabelled example can empty the version space.
\end{tight}
\end{pitfall}

%---------------------------------------------------------------------
\begin{lab}[Candidate Elimination in Forty Lines]
The whole of this chapter, run: \textsc{Find-S}, the boundary trace of the worked
example, and a brute-force check of the version space by listing every
hypothesis.
\begin{lstlisting}[language=Python]
from itertools import product

# Release decisions. "?" = any value, "0" = no value (empty constraint)
VALUES = [("Pass", "Flaky", "Fail"), ("Done", "Pending"),
          ("Complete", "Partial"), ("Ready", "None"), ("Quiet", "Peak"),
          ("Same", "Changed")]
D = [(("Pass", "Done", "Complete", "Ready", "Quiet", "Same"), True),
     (("Pass", "Done", "Partial", "Ready", "Quiet", "Same"), True),
     (("Fail", "Pending", "Partial", "Ready", "Quiet", "Changed"), False),
     (("Pass", "Done", "Partial", "Ready", "Peak", "Changed"), True)]
show = lambda hs: "  ".join("<" + ", ".join(h) + ">" for h in sorted(hs))


def covers(h, x):             # does h classify x as positive?
    return all(c == "?" or c == v for c, v in zip(h, x))


def more_general(h1, h2):     # h1 >=g h2
    return "0" in h2 or all(a == "?" or a == b for a, b in zip(h1, h2))


def generalize(s, x):         # minimal generalisation of s covering x
    return tuple(v if c == "0" else (c if c == v else "?")
                 for c, v in zip(s, x))


def specializations(g, x):    # minimal specialisations of g excluding x
    return [g[:i] + (v,) + g[i + 1:] for i, c in enumerate(g) if c == "?"
            for v in VALUES[i] if v != x[i]]


# --- Find-S --------------------------------------------------------------
h = ("0",) * 6
for x, positive in D:
    if positive:
        h = generalize(h, x)
print("Find-S:", show([h]))

# --- Candidate-Elimination -----------------------------------------------
S, G = {("0",) * 6}, {("?",) * 6}
for k, (x, positive) in enumerate(D, 1):
    if positive:
        G = {g for g in G if covers(g, x)}
        S = {generalize(s, x) for s in S}
        S = {s for s in S if any(more_general(g, s) for g in G)}
    else:
        S = {s for s in S if not covers(s, x)}
        new = {g for g in G if not covers(g, x)}
        for g in G - new:
            new |= {h for h in specializations(g, x)
                    if any(more_general(h, s) for s in S)}
        G = {g for g in new
             if not any(o != g and more_general(o, g) for o in new)}
    print(f"release {k}  S: {show(S)}\n           G: {show(G)}")

# --- brute force: list every hypothesis, keep the consistent ones ---------
H = list(product(*[v + ("?",) for v in VALUES]))
VS = [h for h in H if all(covers(h, x) == p for x, p in D)]
print(len(VS), "hypotheses in the version space:")
print(show(VS))
new_release = ("Pass", "Done", "Complete", "None", "Quiet", "Same")
yes = sum(covers(h, new_release) for h in VS)
print(f"vote on {new_release}: {yes} yes, {len(VS) - yes} no")
\end{lstlisting}
Now add a fifth example, a copy of Release 1 recorded as \emph{No} --- a data-entry
mistake --- and run the trace again. What happens to $S$, $G$ and the version
space?
\end{lab}

%---------------------------------------------------------------------
\begin{checkpoint}
\begin{tightnum}
  \item Is $\langle$Pass, ?, ?, Ready, ?, ?$\rangle \ge_g \langle$Pass, Done,
        ?, Ready, ?, ?$\rangle$? Is the reverse true?
  \item Run \textsc{Find-S} on the release decisions with Release 4 removed. What
        does it return?
  \item Why does an unbiased learner's vote on an unseen instance always end in a
        tie?
\end{tightnum}
\end{checkpoint}

%---------------------------------------------------------------------
\begin{chaptersummary}
\begin{tight}
  \item Concept learning infers a yes/no function from examples by searching a
        hypothesis space --- here, conjunctions of attribute constraints.
  \item Hypotheses are partially ordered from specific to general; the ordering
        lets a learner search without listing every hypothesis.
  \item \textsc{Find-S} returns the most specific hypothesis consistent with the
        positive examples, but cannot say whether it is the only one.
  \item The version space is every hypothesis consistent with the data;
        \textsc{Candidate-Elimination} represents it by its boundaries $S$ and
        $G$, which positive and negative examples move towards each other.
  \item A partially learned concept classifies by vote; the most informative
        query splits the version space in half.
  \item Without an inductive bias there is no generalisation. The bias of
        \textsc{Candidate-Elimination} is that the target is in $H$ --- and it
        assumes noise-free data, which later chapters drop.
\end{tight}
\end{chaptersummary}

\begin{reviewq}
  \item \textbf{(MCQ)} \textsc{Find-S} ignores: (a)~positive examples (b)~negative
        examples (c)~the first example (d)~nothing.
  \item \textbf{(MCQ)} A negative example can change: (a)~$S$ only (b)~$G$ only,
        though it may also remove members of $S$ (c)~neither (d)~the hypothesis
        space.
  \item \textbf{(MCQ)} With 3 attributes of 3 values each, the number of
        semantically distinct conjunctive hypotheses is: (a)~27 (b)~64 (c)~65
        (d)~$2^{27}$.
  \item \textbf{(Short)} Define the version space, and explain why
        \textsc{Candidate-Elimination} does not need to list it.
  \item \textbf{(Short)} State the inductive learning hypothesis. Why is it an
        assumption rather than a theorem?
  \item \textbf{(Short)} Explain how a single mislabelled training example affects
        \textsc{Candidate-Elimination}.
  \item \textbf{(Applied)} Trace \textsc{Candidate-Elimination} on the release
        decisions presented in the order Release 3, Release 1, Release 2,
        Release 4. Is the final version space the same? Is any intermediate
        boundary different?
  \item \textbf{(Applied)} For the final version space of this chapter, find a
        release that splits the vote 2 to 4 other than the one in the worked
        example, and one that splits it 3 to 3. Explain why the second is the
        better query.
\end{reviewq}
